\documentclass[../../main.tex]{subfiles}
\begin{document}

\paragraph{【新課程】}
正$n$角形の頂点を順次$A_1,A_2,\cdots,A_n$とする．
\begin{enumerate}
  \item これらのうちの任意の3点を結んでできる三角形の総数を求めよ．
  \item 上の三角形のうちで鋭角三角形になるものの総数を求めよ．
\end{enumerate}

\paragraph{【旧課程】}
正数列$a_0,a_1,a_2,\cdots,a_n,\cdots\cdots$で条件
\begin{align}
  a_0=1, \quad a_n-a_{n+1}=a_{n+2} \quad (n=0,1,2,\cdots\cdots) \tag{＊}
\end{align}
を満たすものは一組しかないことを，次の順序で証明せよ．
\begin{enumerate}
  \item[(i)] $a_n=\left( \dfrac{\sqrt{5}-1}{2} \right)^n$は条件(＊)を満たす．
  \item[(ii)] 条件(＊)を満たす任意の数列$\{ {a'}_n \}$について，
    数列$c_n=(-1)^n({a'}_n-a_n)$ $(n=0,1,2,\cdots\cdots)$は
    $c_n+c_{n+1}=c_{n+2}$および$|c_n|\geqq(n-1)|c_1|$を満たす．
  \item[(iii)] (ii)の数列$\{ {a'}_n \}$が正数列ならば，${a'}_n\leqq1$，かつ，${a'}_n=a_n$でなければならない．
\end{enumerate}

\end{document}
